% §16 Problem 3, V > v20, at the stationary emission angle (largest theta).
% The tips of v1 and v2 lie on the circles of radii v10 and v20 centred at the
% tip T of V, on opposite sides of T, and theta is the angle at which the
% segment between them is seen from A. The dotted circle has this segment as
% diameter: theta is obtuse when A is inside it.
\documentclass{SolutionFigure}

\begin{document}
\begin{tikzpicture}
    % Parameters (drawing units only)
    \def\va{1}      % v_10
    \def\vb{2}      % v_20
    \def\V{3}     % V
    \pgfmathsetmacro{\th}{acos((\vb-\va)*\V/(\V*\V-\va*\vb))} % emission angle theta_0 (degrees)
    \def\lab{-12}   % direction of the theta label from A (degrees)

    \coordinate (A) at (0,0);
    \coordinate (T) at (\V,0);
    \coordinate (P1) at ($(T)+(\th:\va)$);
    \coordinate (P2) at ($(T)+({\th+180}:\vb)$);
    \coordinate (C) at ($(P1)!0.5!(P2)$);

    % Circles of the C-system speeds, line of emission, circle of diameter P1 P2
    \draw[guide] (T) circle[radius=\va];
    \draw[guide] (T) circle[radius=\vb];
    \draw[thin, dotted] (C) circle[radius={(\va+\vb)/2}];
    \draw[guide] (P2) -- (P1);

    % Velocities
    \draw[vector] (A) -- node[pos=0.5, below=0pt, font=\small] {$\mathbf{V}$} (T);
    \fill (T) circle[radius=0.6pt] node[above right=0pt, font=\scriptsize] {$T$};
    \draw[vector, highlight] (A) -- (P1) node[pos=1, above=0pt] {$\mathbf{v}_1$};
    \draw[vector, highlight] (A) -- (P2) node[pos=1, below=0pt] {$\mathbf{v}_2$};
    \fill (A) circle[radius=0.8pt] node[left=1pt, font=\scriptsize] {$A$};
    \pic[angle arc, angle radius=3mm] {angle = P2--A--P1};
    \node at ($(A)+(\lab:0.75)$) {$\theta$};
\end{tikzpicture}
\end{document}
